\subsection{Memory Access Functions}
\begin{itemize}
\item \forthinlinehl{@}\ \forthinline{( addr -- )}

\noindent Read directly from the memory address and push the value to the TOS

\noindent Arguments: (in order from TOS down, reverse text order):
\begin{enumerate}
\item \forthinline{addr}: (TOS) The address to read from (must be aligned to a 4-byte (32-bit) boundary)
\end{enumerate}

\noindent Return values: none

\item \forthinlinehl{!}\ \forthinline{( x addr -- )}

\noindent Write directly to memory address

\noindent Arguments: (in order from TOS down, reverse text order):
\begin{enumerate}
\item \forthinline{addr}: (TOS) The address to write to (must be aligned to a 4-byte (32-bit) boundary)
\item \forthinline{x}: The value to write to the memory address
\end{enumerate}

\noindent Return values: none

\item \forthinlinehl{+!}\ \forthinline{( x addr -- )}

\noindent Add to value, direct memory access (*addr += x)

\noindent Arguments: (in order from TOS down, reverse text order):
\begin{enumerate}
\item \forthinline{addr}: (TOS) The address to write to (must be aligned to a 4-byte (32-bit) boundary)
\item \forthinline{x}: The value to add to the value at the memory address
\end{enumerate}

\noindent Return values: none

\item \forthinlinehl{sbi}\ \forthinline{( bitnum addr -- )}

\noindent Set bit, direct memory access (*addr {\textbar}= (1 << bitnum))

\noindent Arguments: (in order from TOS down, reverse text order):
\begin{enumerate}
\item \forthinline{addr}: (TOS) The address to write to (must be aligned to a 4-byte (32-bit) boundary)
\item \forthinline{bitnum}: The bit number
\end{enumerate}

\noindent Return values: none

\item \forthinlinehl{cbi}\ \forthinline{( bitnum addr -- )}

\noindent Clear bit, direct memory access (*addr {\&}= {\textasciitilde}(1 << bitnum))

\noindent Arguments: (in order from TOS down, reverse text order):
\begin{enumerate}
\item \forthinline{addr}: (TOS) The address to write to (must be aligned to a 4-byte (32-bit) boundary)
\item \forthinline{bitnum}: The bit number
\end{enumerate}

\noindent Return values: none

\item \forthinlinehl{mask}\ \forthinline{( bitmask x addr -- )}

\noindent Change only masked bits, direct memory access (*addr = (*addr {\&} {\textasciitilde}bitmask) {\textbar} (x {\&} bitmask))

\noindent Arguments: (in order from TOS down, reverse text order):
\begin{enumerate}
\item \forthinline{addr}: (TOS) The address to write to (must be aligned to a 4-byte (32-bit) boundary)
\item \forthinline{x}: The value to write (only masked bits will be written)
\item \forthinline{bitmask}: The bit mask. Only bits with 1s can be changed
\end{enumerate}

\noindent Return values: none

\end{itemize}

\subsection{Print and Input Functions}
\begin{itemize}

\item \forthinlinehl{.}\ \forthinline{( x -- )}

\noindent Prints the TOS as an integer

\noindent Arguments: (in order from TOS down, reverse text order):
\begin{enumerate}
\item \forthinline{x}: (TOS) The value to print
\end{enumerate}

\noindent Return values: none

\item \forthinlinehl{h.}\ \forthinline{( x -- )}

\noindent Prints the TOS as a hex string

\noindent Arguments: (in order from TOS down, reverse text order):
\begin{enumerate}
\item \forthinline{x}: (TOS) The value to print
\end{enumerate}

\noindent Return values: none

\item \forthinlinehl{echo}\ \forthinline{( bool -- )}

\noindent Change the echo state, which controls whether the Forth interpreter repeats all of its received characters

\noindent Arguments: (in order from TOS down, reverse text order):
\begin{enumerate}
\item \forthinline{bool}: (TOS) If 0, does not echo. Otherwise, it does
\end{enumerate}

\noindent Return values: none

\item \forthinlinehl{emit}\ \forthinline{( x -- )}

\noindent Prints the char at the TOS (integers are mapped to their ASCII counterparts)

\noindent Arguments: (in order from TOS down, reverse text order):
\begin{enumerate}
\item \forthinline{x}: (TOS) The ASCII value of the char to print
\end{enumerate}

\noindent Return values: none

\item \forthinlinehl{key}\ \forthinline{( -- char )}

\noindent Waits for a char over UART and pushes the ASCII value of the char to the TOS

\noindent Arguments: none

\noindent Return values: (in order from TOS down, reverse text order):
\begin{enumerate}
\item \forthinline{char}: (TOS) The char that was received over UART
\end{enumerate}

\item \forthinlinehl{list}\ \forthinline{( -- )}

\noindent Prints all currently available functions

\noindent Arguments: none

\noindent Return values: none

\end{itemize}

\subsection{Arithmetic Functions}
\begin{itemize}

\item \forthinlinehl{+}\ \forthinline{( y x -- ret )}

\noindent Adds two numbers (y + x)

\noindent Arguments: (in order from TOS down, reverse text order):
\begin{enumerate}
\item \forthinline{x}: (TOS) Signed integer
\item \forthinline{y}: Signed integer
\end{enumerate}

\noindent Return values: (in order from TOS down, reverse text order):
\begin{enumerate}
\item \forthinline{ret}: (TOS) Return value
\end{enumerate}

\item \forthinlinehl{-}\ \forthinline{( y x -- ret )}

\noindent Subtracts two numbers (y - x)

\noindent Arguments: (in order from TOS down, reverse text order):
\begin{enumerate}
\item \forthinline{x}: (TOS) Signed integer
\item \forthinline{y}: Signed integer
\end{enumerate}

\noindent Return values: (in order from TOS down, reverse text order):
\begin{enumerate}
\item \forthinline{ret}: (TOS) Return value
\end{enumerate}

\item \forthinlinehl{*}\ \forthinline{( y x -- retprodhi retprodlo )}

\noindent Multiplies two numbers (y * x)

\noindent Arguments: (in order from TOS down, reverse text order):
\begin{enumerate}
\item \forthinline{x}: (TOS) Signed integer
\item \forthinline{y}: Signed integer
\end{enumerate}

\noindent Return values: (in order from TOS down, reverse text order):
\begin{enumerate}
\item \forthinline{retprodlo}: (TOS) Lower 32 bits of the return product
\item \forthinline{retprodhi}: Upper 32 bits of the return product
\end{enumerate}

\item \forthinlinehl{/{\%}}\ \forthinline{( y x -- retdiv retrem )}

\noindent Divides two numbers with remainder (y / x and y {\%} x)

\noindent Arguments: (in order from TOS down, reverse text order):
\begin{enumerate}
\item \forthinline{x}: (TOS) Signed integer
\item \forthinline{y}: Signed integer
\end{enumerate}

\noindent Return values: (in order from TOS down, reverse text order):
\begin{enumerate}
\item \forthinline{retrem}: (TOS) The remainder (or modulo) result
\item \forthinline{retdiv}: The integer division result
\end{enumerate}

\item \forthinlinehl{neg}\ \forthinline{( x -- ret )}

\noindent Returns the negative (twos complement) of x

\noindent Arguments: (in order from TOS down, reverse text order):
\begin{enumerate}
\item \forthinline{x}: (TOS) Signed integer
\end{enumerate}

\noindent Return values: (in order from TOS down, reverse text order):
\begin{enumerate}
\item \forthinline{ret}: (TOS) Return value
\end{enumerate}

\item \forthinlinehl{abs}\ \forthinline{( x -- ret )}

\noindent Absolute value ({\textbar}x{\textbar})

\noindent Arguments: (in order from TOS down, reverse text order):
\begin{enumerate}
\item \forthinline{x}: (TOS) Signed integer
\end{enumerate}

\noindent Return values: (in order from TOS down, reverse text order):
\begin{enumerate}
\item \forthinline{ret}: (TOS) Return value
\end{enumerate}

\end{itemize}

\subsection{Bitwise Logic Functions}
\begin{itemize}

\item \forthinlinehl{{\~}}\ \forthinline{( x -- ret )}

\noindent Bitwise complement/inversion ({\textasciitilde}x)

\noindent Arguments: (in order from TOS down, reverse text order):
\begin{enumerate}
\item \forthinline{x}: (TOS) Signed integer
\end{enumerate}

\noindent Return values: (in order from TOS down, reverse text order):
\begin{enumerate}
\item \forthinline{ret}: (TOS) Return value
\end{enumerate}

\item \forthinlinehl{|}\ \forthinline{( y x -- ret )}

\noindent Bitwise OR (y {\textbar} x)

\noindent Arguments: (in order from TOS down, reverse text order):
\begin{enumerate}
\item \forthinline{x}: (TOS) Signed integer
\item \forthinline{y}: Signed integer
\end{enumerate}

\noindent Return values: (in order from TOS down, reverse text order):
\begin{enumerate}
\item \forthinline{ret}: (TOS) Return value
\end{enumerate}

\item \forthinlinehl{&}\ \forthinline{( y x -- ret )}

\noindent Bitwise AND (y {\&} x)

\noindent Arguments: (in order from TOS down, reverse text order):
\begin{enumerate}
\item \forthinline{x}: (TOS) Signed integer
\item \forthinline{y}: Signed integer
\end{enumerate}

\noindent Return values: (in order from TOS down, reverse text order):
\begin{enumerate}
\item \forthinline{ret}: (TOS) Return value
\end{enumerate}

\item \forthinlinehl{^}\ \forthinline{( y x -- ret )}

\noindent Bitwise XOR (y {\textasciicircum} x)

\noindent Arguments: (in order from TOS down, reverse text order):
\begin{enumerate}
\item \forthinline{x}: (TOS) Signed integer
\item \forthinline{y}: Signed integer
\end{enumerate}

\noindent Return values: (in order from TOS down, reverse text order):
\begin{enumerate}
\item \forthinline{ret}: (TOS) Return value
\end{enumerate}

\item \forthinlinehl{*2}\ \forthinline{( x -- ret )}

\noindent Shift left 1 bit (x << 1)

\noindent Arguments: (in order from TOS down, reverse text order):
\begin{enumerate}
\item \forthinline{x}: (TOS) Signed integer
\end{enumerate}

\noindent Return values: (in order from TOS down, reverse text order):
\begin{enumerate}
\item \forthinline{ret}: (TOS) Return value
\end{enumerate}

\item \forthinlinehl{/2}\ \forthinline{( x -- ret )}

\noindent Shift right 1 bit (x >> 1)

\noindent Arguments: (in order from TOS down, reverse text order):
\begin{enumerate}
\item \forthinline{x}: (TOS) Signed integer
\end{enumerate}

\noindent Return values: (in order from TOS down, reverse text order):
\begin{enumerate}
\item \forthinline{ret}: (TOS) Return value
\end{enumerate}

\item \forthinlinehl{sll}\ \forthinline{( x bits -- ret )}

\noindent Shift left logical (x << bits)

\noindent Arguments: (in order from TOS down, reverse text order):
\begin{enumerate}
\item \forthinline{bits}: (TOS) Signed integer
\item \forthinline{x}: Signed integer
\end{enumerate}

\noindent Return values: (in order from TOS down, reverse text order):
\begin{enumerate}
\item \forthinline{ret}: (TOS) Return value
\end{enumerate}

\item \forthinlinehl{srl}\ \forthinline{( x bits -- ret )}

\noindent Shift right logical (x >> bits), no sign extension

\noindent Arguments: (in order from TOS down, reverse text order):
\begin{enumerate}
\item \forthinline{bits}: (TOS) Signed integer
\item \forthinline{x}: Signed integer
\end{enumerate}

\noindent Return values: (in order from TOS down, reverse text order):
\begin{enumerate}
\item \forthinline{ret}: (TOS) Return value
\end{enumerate}

\end{itemize}

\subsection{Comparison Functions}
\begin{itemize}

\item \forthinlinehl{>}\ \forthinline{( y x -- ret )}

\noindent Greater than. Returns 1 if y > x, returns 0 otherwise.

\noindent Arguments: (in order from TOS down, reverse text order):
\begin{enumerate}
\item \forthinline{x}: (TOS) Signed integer
\item \forthinline{y}: Signed integer
\end{enumerate}

\noindent Return values: (in order from TOS down, reverse text order):
\begin{enumerate}
\item \forthinline{ret}: (TOS) 1 if true, 0 if false
\end{enumerate}

\item \forthinlinehl{<}\ \forthinline{( y x -- ret )}

\noindent Less than. Returns 1 if y < x, returns 0 otherwise.

\noindent Arguments: (in order from TOS down, reverse text order):
\begin{enumerate}
\item \forthinline{x}: (TOS) Signed integer
\item \forthinline{y}: Signed integer
\end{enumerate}

\noindent Return values: (in order from TOS down, reverse text order):
\begin{enumerate}
\item \forthinline{ret}: (TOS) 1 if true, 0 if false
\end{enumerate}

\item \forthinlinehl{==}\ \forthinline{( y x -- ret )}

\noindent Equals. Returns 1 if y == x, returns 0 otherwise.

\noindent Arguments: (in order from TOS down, reverse text order):
\begin{enumerate}
\item \forthinline{x}: (TOS) Signed integer
\item \forthinline{y}: Signed integer
\end{enumerate}

\noindent Return values: (in order from TOS down, reverse text order):
\begin{enumerate}
\item \forthinline{ret}: (TOS) 1 if true, 0 if false
\end{enumerate}

\end{itemize}
\subsection{Boolean Functions}
\begin{itemize}

\item \forthinlinehl{not}\ \forthinline{( x -- ret )}

\noindent Logical not (!x). Returns 1 if x is equal to 0, returns 0 otherwise

\noindent Arguments: (in order from TOS down, reverse text order):
\begin{enumerate}
\item \forthinline{x}: (TOS) Signed integer
\end{enumerate}

\noindent Return values: (in order from TOS down, reverse text order):
\begin{enumerate}
\item \forthinline{ret}: (TOS) Return value
\end{enumerate}

\item \forthinlinehl{and}\ \forthinline{( y x -- ret )}

\noindent Logical and (y {\&}{\&} x). Returns 1 if y and x are both true (nonzero), returns 0 otherwise

\noindent Arguments: (in order from TOS down, reverse text order):
\begin{enumerate}
\item \forthinline{x}: (TOS) Signed integer
\item \forthinline{y}: Signed integer
\end{enumerate}

\noindent Return values: (in order from TOS down, reverse text order):
\begin{enumerate}
\item \forthinline{ret}: (TOS) Return value
\end{enumerate}

\item \forthinlinehl{or}\ \forthinline{( y x -- ret )}

\noindent Logical or (y {\textbar}{\textbar} x). Returns 1 if y or x are true (nonzero), returns 0 otherwise

\noindent Arguments: (in order from TOS down, reverse text order):
\begin{enumerate}
\item \forthinline{x}: (TOS) Signed integer
\item \forthinline{y}: Signed integer
\end{enumerate}

\noindent Return values: (in order from TOS down, reverse text order):
\begin{enumerate}
\item \forthinline{ret}: (TOS) Return value
\end{enumerate}

\end{itemize}
\subsection{Stack Manipulation Functions}
\begin{itemize}

\item \forthinlinehl{dup}\ \forthinline{( x -- x x )}

\noindent Duplicates the TOS (pops the TOS and then pushes the value to the TOS twice)

\noindent Arguments: (in order from TOS down, reverse text order):
\begin{enumerate}
\item \forthinline{x}: (TOS) Value to duplicate
\end{enumerate}

\noindent Return values: (in order from TOS down, reverse text order):
\begin{enumerate}
\item \forthinline{x}: (TOS) The duplicated value
\item \forthinline{x}: The original TOS
\end{enumerate}

\item \forthinlinehl{drop}\ \forthinline{( x -- )}

\noindent Pops the TOS and sets it as the internal Forth \forthinline{i} value. Usually used to delete the TOS.

\noindent Arguments: (in order from TOS down, reverse text order):
\begin{enumerate}
\item \forthinline{x}: (TOS) The value that is dropped
\end{enumerate}

\noindent Return values: none

\item \forthinlinehl{swap}\ \forthinline{( y x -- x y )}

\noindent Swaps the TOS with the value below the TOS

\noindent Arguments: (in order from TOS down, reverse text order):
\begin{enumerate}
\item \forthinline{x}: (TOS) Old TOS
\item \forthinline{y}: Old value below the TOS
\end{enumerate}

\noindent Return values: (in order from TOS down, reverse text order):
\begin{enumerate}
\item \forthinline{y}: (TOS) Old value below the TOS
\item \forthinline{x}: Old TOS
\end{enumerate}

\item \forthinlinehl{over}\ \forthinline{( y x -- y x y )}

\noindent Copies the value below the TOS and pushes it to the TOS

\noindent Arguments: (in order from TOS down, reverse text order):
\begin{enumerate}
\item \forthinline{x}: (TOS) Signed Integer
\item \forthinline{y}: Signed Integer
\end{enumerate}

\noindent Return values: (in order from TOS down, reverse text order):
\begin{enumerate}
\item \forthinline{y}: (TOS) Signed Integer
\item \forthinline{x}: Signed Integer
\item \forthinline{y}: Signed Integer
\end{enumerate}

\item \forthinlinehl{tuck}\ \forthinline{( y x -- x y x )}

\noindent Copies the TOS and inserts it two after the TOS

\noindent Arguments: (in order from TOS down, reverse text order):
\begin{enumerate}
\item \forthinline{x}: (TOS) Signed Integer
\item \forthinline{y}: Signed Integer
\end{enumerate}

\noindent Return values: (in order from TOS down, reverse text order):
\begin{enumerate}
\item \forthinline{x}: (TOS) Signed Integer
\item \forthinline{y}: Signed Integer
\item \forthinline{x}: Signed Integer
\end{enumerate}

\item \forthinlinehl{roll}\ \forthinline{( n -- )}

\noindent Removes value at index n of the stack and pushes it to the TOS (zero indexed)

\noindent Arguments: (in order from TOS down, reverse text order):
\begin{enumerate}
\item \forthinline{n}: (TOS) Stack index (zero indexed)
\end{enumerate}

\noindent Return values: none

\item \forthinlinehl{pick}\ \forthinline{( n -- )}

\noindent Copies value at index n of the stack and pushes it to the TOS (zero indexed)

\noindent Arguments: (in order from TOS down, reverse text order):
\begin{enumerate}
\item \forthinline{n}: (TOS) Stack index (zero indexed)
\end{enumerate}

\noindent Return values: none

\item \forthinlinehl{depth}\ \forthinline{( -- ret )}

\noindent Gets size of stack and pushes it to the top of the stack (making the new stack size the value of TOS + 1)

\noindent Arguments: none

\noindent Return values: (in order from TOS down, reverse text order):
\begin{enumerate}
\item \forthinline{ret}: (TOS) Return value
\end{enumerate}

\end{itemize}

\subsection{SPI Flash R/W and Programming Functions} \label{ss:forthflash}
\begin{itemize}

\item \forthinlinehl{fr}\ \forthinline{( bin length addr -- )}

\noindent Reads data from the SPI flash.

\noindent Arguments: (in order from TOS down, reverse text order):
\begin{enumerate}
\item \forthinline{addr}: (TOS) The address in the SPI flash to begin reading from
\item \forthinline{length}: The number of bytes to read
\item \forthinline{bin}: If true (nonzero), transmits binary data, otherwise transmits ASCII hex string
\end{enumerate}

\noindent Return values: none

\item \forthinlinehl{fe}\ \forthinline{( conf addr -- eraseSuccessful )}

\noindent Erases a 256-byte page from the SPI flash memory. The start address must be 256-byte aligned.

\noindent Arguments: (in order from TOS down, reverse text order):
\begin{enumerate}
\item \forthinline{addr}: (TOS) The start address of the page to erase in the SPI flash. Must be 256-byte aligned.
\item \forthinline{conf}: If equal to 123, erases the page and returns 1, otherwise does not erase and returns 0
\end{enumerate}

\noindent Return values: (in order from TOS down, reverse text order):
\begin{enumerate}
\item \forthinline{eraseSuccessful}: (TOS) Nonzero if successful, zero if failed
\end{enumerate}

\item \forthinlinehl{fw}\ \forthinline{( bin addr -- )}

\noindent Writes data to the SPI flash. Must write exactly 256 bytes at a time. The start address must be 256-byte aligned.

\noindent Arguments: (in order from TOS down, reverse text order):
\begin{enumerate}
\item \forthinline{addr}: (TOS) The address in the SPI flash to begin writing to. Must be 256-byte aligned.
\item \forthinline{bin}: If true (nonzero), expects to receive binary data, otherwise expects to receive an ASCII hex string
\end{enumerate}

\noindent Return values: none

\end{itemize}

\subsection{Miscellaneous Functions}
\begin{itemize}

\item \forthinlinehl{rst}\ \forthinline{( -- )}

\noindent Performs a soft reset (turns on all memory cells and jumps to address 0 in the ROM, does not reset clocks or hardware)

\noindent Arguments: none

\noindent Return values: none

\item \forthinlinehl{clk}\ \forthinline{( meastime clksel -- freq )}

\noindent Measures selected clock frequency

\noindent Arguments: (in order from TOS down, reverse text order):
\begin{enumerate}
\item \forthinline{clksel}: (TOS) Clock select. 0 = SMCLK; 1 = MCLK; 2 = LFXT; 3 = HFXT
\item \forthinline{meastime}: Measurement time. 0 = 1 s; 1 = 0.5 s; 2 = 0.25 s; 3 = 0.125 s
\end{enumerate}

\noindent Return values: (in order from TOS down, reverse text order):
\begin{enumerate}
\item \forthinline{freq}: (TOS) Selected clock frequency in Hz
\end{enumerate}

\item \forthinlinehl{min}\ \forthinline{( y x -- ret )}

\noindent Returns minimum value of y and x

\noindent Arguments: (in order from TOS down, reverse text order):
\begin{enumerate}
\item \forthinline{x}: (TOS) Signed integer
\item \forthinline{y}: Signed integer
\end{enumerate}

\noindent Return values: (in order from TOS down, reverse text order):
\begin{enumerate}
\item \forthinline{ret}: (TOS) Return value
\end{enumerate}

\item \forthinlinehl{max}\ \forthinline{( y x -- ret )}

\noindent Returns maximum value of y and x

\noindent Arguments: (in order from TOS down, reverse text order):
\begin{enumerate}
\item \forthinline{x}: (TOS) Signed integer
\item \forthinline{y}: Signed integer
\end{enumerate}

\noindent Return values: (in order from TOS down, reverse text order):
\begin{enumerate}
\item \forthinline{ret}: (TOS) Return value
\end{enumerate}

\item \forthinlinehl{swpb}\ \forthinline{( x -- ret )}

\noindent Swap bits 15 downto 8 with bits 7 downto 0 of TOS, also causes bits 31 downto 16 to all be 0

\noindent Arguments: (in order from TOS down, reverse text order):
\begin{enumerate}
\item \forthinline{x}: (TOS) Signed integer
\end{enumerate}

\noindent Return values: (in order from TOS down, reverse text order):
\begin{enumerate}
\item \forthinline{ret}: (TOS) Return value
\end{enumerate}

\item \forthinlinehl{swphw}\ \forthinline{( x -- ret )}

\noindent Swap upper 16 bits with lower 16 bits of TOS

\noindent Arguments: (in order from TOS down, reverse text order):
\begin{enumerate}
\item \forthinline{x}: (TOS) Signed integer
\end{enumerate}

\noindent Return values: (in order from TOS down, reverse text order):
\begin{enumerate}
\item \forthinline{ret}: (TOS) Return value
\end{enumerate}

\end{itemize}

